\section{Uniform factorial convolution and the hyperbolic amplitudes}
\label{sec:leading}

The source conjecture correctly identifies the contribution of any fixed
number of factors in \eqref{eq:C-finite}. The missing step is to control
that number uniformly. A factorial convolution estimate supplies this
control without fitting late coefficients or assuming resurgence.

\subsection{A majorant that retains the endpoint factorial}

Scale the coefficients by
\begin{equation}\label{eq:positive-scaled}
 a_m=A^m(-1)^mr_m,\qquad a_1=c,\qquad
 a_{2j}=4b_{2j}\zeta(2j)^2,\qquad a_{2j+1}=0\ (j\ge1),
\end{equation}
and set
\begin{equation}\label{eq:S-D}
 S_{n,k}=[x^n]\left(\sum_{m\ge1}a_mx^m\right)^k,
 \qquad D_n(u)=(-1)^{n+1}A^nC_n(u).
\end{equation}
All $a_m$ are nonnegative. If
$K\ge\max\{c,4\zeta(2)^2\}$, then $a_m\le Kb_m$ for every $m$.
Equation~\eqref{eq:C-finite} becomes
\begin{equation}\label{eq:D-formula}
 D_n(u)=\sum_{k=1}^n\frac{u^k}{k!}S_{n,k}T_{n,k}(u).
\end{equation}

\begin{lemma}[Endpoint-preserving factorial bounds]\label{lem:factorial}
For $n\ge2$,
\begin{equation}\label{eq:b-convolution}
 \sum_{m=1}^{n-1}b_mb_{n-m}\le3b_{n-1}.
\end{equation}
Consequently, for $1\le k\le n$,
\begin{equation}\label{eq:S-bound}
 S_{n,k}\le K^k3^{k-1}b_{n-k+1},\qquad
 \frac{n^{k-1}b_{n-k+1}}{b_n}\le ke^k.
\end{equation}
\end{lemma}
\begin{proof}
The ratio of the $m$th summand in \eqref{eq:b-convolution} to $b_{n-1}$ is
\[
 \frac{n-1}{m(n-m)}\binom{n-2}{m-1}^{-1}
 \le\binom{n-2}{m-1}^{-1}.
\]
For $N\ge2$ every interior binomial coefficient $\binom Nj$ is at least
$N$, so
$\sum_{j=0}^N\binom Nj^{-1}\le2+(N-1)/N<3$.
The cases $N=0,1$ satisfy the same bound. This proves
\eqref{eq:b-convolution}. Induction on $k$, shifting the convolution index
at each step, gives
$[x^n](\sum_{m\ge1}b_mx^m)^k\le3^{k-1}b_{n-k+1}$, and therefore the first
part of \eqref{eq:S-bound}.

For the second part let $h=k-1$. Direct factorial cancellation gives
\begin{equation}\label{eq:ratio-exact}
 \frac{n^hb_{n-h}}{b_n}
 =\frac{n}{n-h}\prod_{i=1}^h\left(1-\frac in\right)^{-1}.
\end{equation}
The first factor is at most $k$, since $(k-1)(n-k)\ge0$.
For the increasing function $f(x)=-\log(1-x)$,
\[
 \sum_{i=1}^{k-1}f(i/n)
 \le n\int_0^{k/n}f(x)\dd x
 =k+(n-k)\log(1-k/n)\le k.
\]
At $k=n$ the integral is finite and $0\log0$ is interpreted as zero.
Exponentiation proves the bound, including the endpoint $k=n$.
\end{proof}

For $|u|\le U$, the exact polynomial \eqref{eq:T} satisfies
\begin{equation}\label{eq:T-bound}
 |T_{n,k}(u)|\le n^{k-1}(1+U)^{k-1},
\end{equation}
because $(k)_j\le(2n)^j$. Combining this with Lemma~\ref{lem:factorial}
gives the central estimate
\begin{equation}\label{eq:master-majorant}
 \frac1{b_n}\left|\frac{u^k}{k!}S_{n,k}T_{n,k}(u)\right|
 \le \frac{k}{k!}K^k3^{k-1}e^kU^k(1+U)^{k-1}
 \le \frac{kB_U^k}{k!}
\end{equation}
for a suitable fixed $B_U$. The final majorant is summable in $k$ and
independent of $n$. It removes the spurious $e^{O(\sqrt n)}$ loss of the
coarser composition count in the predecessor.

The same conclusion holds after division by $u$ for even $n$, and by
$u^2$ for odd $n\ge3$. In the latter case the exact sum begins at $k=2$,
by \eqref{eq:valuation}. On a compact disk one uses
$|u|^{k-v}\le\max(1,U)^k$, with $v=1$ or $2$. Every quotient is a
polynomial, so no estimate divides by a small nonzero number.

\subsection{The contribution of one large index}

\begin{lemma}[First endpoint extraction]\label{lem:endpoint}
For $k\ge2$ and $n\ge k+2$,
\begin{equation}\label{eq:endpoint}
 S_{n,k}=k c^{k-1}a_{n-k+1}+E_{n,k},\qquad
 0\le E_{n,k}\le6\binom{k}{2}K^k3^{k-2}b_{n-k}.
\end{equation}
For $k=1$ the main term is exact.
\end{lemma}
\begin{proof}
Compositions with exactly one part greater than $1$ give the main term.
Every remaining composition has two parts at least $2$. Mark a pair of
such positions and use a union bound over $\binom{k}{2}$ choices.

For $N\ge4$,
\begin{equation}\label{eq:interior-b}
 \sum_{m=2}^{N-2}b_mb_{N-m}\le6b_{N-2}.
\end{equation}
Indeed, put $M=N-4$ and $j=m-2$. The ratio of a summand to $b_{N-2}$ is
\[
 \frac{M+2}{M+1}
 \frac{(j+1)(M-j+1)}{(j+2)(M-j+2)}\binom Mj^{-1}
 \le2\binom Mj^{-1}.
\]
The reciprocal-binomial bound proves \eqref{eq:interior-b}.
After summing the two marked parts, convolve with the remaining $k-2$
parts and enlarge the index range. The result is bounded by
\[
 6K^k[x^{n-2}]\left(\sum_{m\ge1}b_mx^m\right)^{k-1}
 \le6K^k3^{k-2}b_{n-k}.
\]
Summing over the marked pair proves \eqref{eq:endpoint}.
\end{proof}

For fixed $k$, multiplication of the remainder in \eqref{eq:endpoint} by
$n^{k-1}/b_n$ gives $O_k(n^{-1})$. The main index is $m=n-k+1$.
If $m$ is even, then
$n^{k-1}a_m/b_n\to4$; if $m$ is odd it is eventually zero.
Also $T_{n,k}(u)/n^{k-1}\to1$ locally uniformly in $u$.
Thus the limit of the $k$th summand in \eqref{eq:D-formula}, divided by
$b_n$, is
\begin{equation}\label{eq:fixed-k-limit}
 \frac{4k u^kc^{k-1}}{k!}
 =4u\frac{(cu)^{k-1}}{(k-1)!}
\end{equation}
when $n-k+1$ is even, and zero in the other parity.
The majorant \eqref{eq:master-majorant} now justifies summation of these
limits. Even $n$ selects even $k-1$ and produces $\cosh(cu)$; odd $n$
selects odd $k-1$ and produces $\sinh(cu)$.

\subsection{Uniform errors, including the zero-parameter boundary}

\begin{theorem}[A quantified proof of the source conjecture]
\label{thm:leading-rate}
For every $U>0$, uniformly for complex $|u|\le U$,
\begin{align}
 \frac{A^nC_n(u)}{4b_nu}&=-\cosh(cu)+O_U(n^{-1})
       &&(n\text{ even}),\label{eq:leading-even}\\
 \frac{A^nC_n(u)}{4b_nu^2}&=\frac{\sinh(cu)}u+O_U(n^{-1})
       &&(n\ge3\text{ odd}).\label{eq:leading-odd}
\end{align}
All quotients at $u=0$ are interpreted by continuity.
\end{theorem}
\begin{proof}
Split the finite sum at $k=n/2$ and first take $k\le n/2$. Let
$h=k-1$ and $R_{n,k}=n^hb_{n-h}/b_n$. The logarithm of
\eqref{eq:ratio-exact} and $-\log(1-x)\le2x$ on $[0,1/2]$ give
\begin{align}
 0\le\log R_{n,k}&\le\frac{h(h+3)}n,\nonumber\\
 |R_{n,k}-1|&\le\frac{ke^kh(h+3)}n,\label{eq:ratio-rate}\\
 \frac{n^hb_{n-h-1}}{b_n}
 &=R_{n,k}\frac{n-h}{(n-h-1)^2}
 \le\frac{4ke^k}{n}\quad(n\ge4).\nonumber
\end{align}
The last estimate controls \eqref{eq:endpoint}. Using the sharper
$(k)_j\le(2k)^j$ in \eqref{eq:T} gives
\begin{equation}\label{eq:T-rate}
 \left|\frac{T_{n,k}(u)}{n^h}-1\right|
 \le(1+kU/n)^h-1
 \le\frac{k^2U}{n}(1+U)^k.
\end{equation}
For the even endpoint $m=n-k+1\ge n/2$, the error
$\zeta(m)^2-1=O(2^{-m})$ is uniformly $O(n^{-1})$.
Combining these estimates bounds the difference between each normalized
summand and \eqref{eq:fixed-k-limit} by
\[
 \frac1n\frac{P(k)\widetilde B_U^k}{k!}
\]
for a fixed polynomial $P$. The bound also holds after division by the
valuation factor $u$ or $u^2$, as explained after
\eqref{eq:master-majorant}. It is summable over $k$.

For $k>n/2$ both the exact summands and the limiting exponential summands
have a majorant of the form $kB_U^k/k!$. The inequality $1\le2k/n$ makes
the summed tail $O_U(n^{-1})$. More generally, using
$1\le(2k/n)^M$ makes this tail $O_{U,M}(n^{-M})$ for any fixed $M$.
Summing \eqref{eq:fixed-k-limit} with its parity restriction proves
\eqref{eq:leading-even}--\eqref{eq:leading-odd}.
\end{proof}

For real $0\le u\le U$, one has $\cosh(cu)\ge1$ and
$\sinh(cu)/u\ge c$, including the limiting value at zero. Dividing by
these lower bounds proves the uniform relative versions
\eqref{eq:even-main}--\eqref{eq:odd-main}. This is stronger than an estimate
whose constant depends on a fixed lower bound for $u$.

\begin{corollary}[Two limits may be taken together]\label{cor:moving-u}
Let $u_n>0$, $u_n\to0$, and $n\to\infty$ through the indicated parity.
Then
\begin{equation}\label{eq:small-u-late}
 C_n(u_n)\sim
 \begin{cases}
 -4u_n b_nA^{-n},&n\text{ even},\\
 A u_n^2 b_nA^{-n},&n\text{ odd}.
 \end{cases}
\end{equation}
The rate at which $u_n$ approaches zero is unrestricted.
\end{corollary}
\begin{proof}
Apply the uniform relative bounds and use
$\cosh(cu)=1+O(u^2)$ and $\sinh(cu)=cu\{1+O(u^2)\}$.
\end{proof}

The coefficient theorem therefore applies to the physical slow parameter
$u=1/X$ while the sector index increases. It does not assert that the
divergent infinite sector series may be summed ordinarily at that physical
pair, nor does a coefficient estimate alone determine an optimally
truncated remainder of the actual inverse.

\section{Finite phases and changes of the initial coefficients}\label{sec:phases}

The hyperbolic functions reflect the two phases $1$ and $-1$ in the
forward coefficients. The same proof yields a finite-phase transformation
law, with no positivity requirement.

\begin{theorem}[Leading amplitudes for finitely many phases]
\label{thm:finite-phases}
Let $A>0$, $b_n=\Gamma(n)/n$, and
$\widetilde R(t)=\sum_{n\ge1}\widetilde r_nt^n$, with
$\widetilde r_1=\rho\in\C$. Suppose that for a finite set
$E\subset\{\omega\in\C:|\omega|=1\}$ of distinct phases,
\begin{equation}\label{eq:phase-input}
 |\widetilde r_n|\le K A^{-n}b_n,\qquad
 \frac{A^n\widetilde r_n}{b_n}
 =\sum_{\omega\in E}\kappa_\omega\omega^n+O(n^{-1}).
\end{equation}
Let $\widetilde C_n(u)$ be the coefficients of the zero-constant formal
solution of \eqref{eq:implicit} with $R$ replaced by $\widetilde R$.
Then, locally uniformly for $u\in\C$,
\begin{equation}\label{eq:phase-output}
 \widetilde C_n(u)=-u A^{-n}b_n
 \left\{\sum_{\omega\in E}\kappa_\omega\omega^n
                   e^{A\rho u/\omega}+O_U(n^{-1})\right\}.
\end{equation}
The bracketed formulation extends through $u=0$. Its error corresponds to
an absolute $O_U(|u|A^{-n}b_n/n)$ error in $\widetilde C_n$.
\end{theorem}
\begin{proof}
Apply the exact formula \eqref{eq:C-finite} to the new input and scale by
$A^n$. The first endpoint is now
$k(A\rho)^{k-1}(A^{n-k+1}\widetilde r_{n-k+1})$.
All remainder estimates of Lemmas~\ref{lem:factorial} and
\ref{lem:endpoint} apply to absolute values: the positivity assumption
was used to construct a majorant, not to establish an identity.
For $k\le n/2$, the remainder in \eqref{eq:phase-input} is $O(n^{-1})$
at the endpoint. The same proof as Theorem~\ref{thm:leading-rate} therefore
justifies the uniform $k$ summation with error $O_U(n^{-1})$.
For each phase the limiting sum is
\[
 -\sum_{k\ge1}\frac{u^k}{(k-1)!}(A\rho)^{k-1}
                 \kappa_\omega\omega^{n-k+1}
 =-u\kappa_\omega\omega^n e^{A\rho u/\omega}.
\]
The sum over $E$ is finite, which proves \eqref{eq:phase-output}.
\end{proof}

The action corresponding to $\omega$ is $\lambda_\omega=A/\omega$.
Its leading amplitude is multiplied by $e^{\lambda_\omega\rho u}$.
This multiplier comes entirely from repeated linear factors of the
forward series. The terms of $T_{n,k}$ with $j\ge1$, which encode the
quadratic part of the inverse core, enter one inverse-index order later.
Cancellation between distinct phases can occur on subsequences, so
\eqref{eq:phase-output} is an additive expansion. A relative equivalent
requires a separate nonvanishing argument.

\begin{corollary}[Polynomial changes of the forward tail]
\label{cor:prefix}
For the Lambert input, replace $R(t)$ by $R(t)+P(t)$, where $P(0)=0$ is a
fixed polynomial and $p_1=[t]P(t)$. The late coefficient formula is
\begin{equation}\label{eq:prefix}
 \widetilde C_n(u)=-2uA^{-n}b_n
 \left\{e^{A(p_1-1/4)u}+(-1)^n e^{-A(p_1-1/4)u}
                     +O_U(n^{-1})\right\}.
\end{equation}
In particular, changes satisfying $P(t)=O(t^2)$ leave both leading action
amplitudes unchanged.
\end{corollary}
\begin{proof}
A polynomial changes no sufficiently late forward coefficient. Apply
Theorem~\ref{thm:finite-phases} with $E=\{1,-1\}$ and
$\kappa_1=\kappa_{-1}=2$.
\end{proof}

This corollary gives a useful separation of data: the arithmetic tail
locates the leading actions, the linear term changes their amplitudes,
and the quadratic term first appears in the correction computed next.
